%!TEX root = main.tex
    \documentclass[
        12pt,
        oneside,
        titlepage,
        final,
        fleqn,
    ]{\myClass}
\usepackage[
        papersize={166.5mm,222mm}, %trail on error zodat er net zoveel tekst op een regel past als bij A4 versie
        top=14mm,
        head=10mm,
        headsep=5mm,
        bottom=15mm,
        footskip=11.7mm,
        left=10mm,
        right=10mm,
        twoside=false
    ]{geometry}

\usepackage{etoolbox}
\usepackage[utf8]{inputenc}
\usepackage[T1]{fontenc}
\usepackage{setspace}
\usepackage{mathptmx}
\usepackage[scaled=.92]{helvet}
\usepackage{charter}
\usepackage[charter]{mathdesign} 
\usepackage[scaled]{berasans} %ifcharacter
\usepackage{cmbright}
\usepackage[default,scale=0.95]{opensans}
\usepackage[varqu]{zi4}
\usepackage[hyphens]{url}
\usepackage{amsmath}
\usepackage{mathtools}
\usepackage{expl3}
\usepackage{xstring}

\newtoggle{onesided}
\ifdef{\forceOneSided}{\toggletrue{onesided}}{\togglefalse{onesided}}

\ifdef{\titel}{}{\newcommand{\titel}{}}
\ifdef{\auteurs}{}{\newcommand{\auteurs}{}}
\ifdef{\documenttype}{}{\newcommand{\documenttype}{}}

\setlength{\jot}{\baselineskip}
%>extramathcommands
\renewcommand{\vec}[1]{\mathbf{#1}}
\newcommand{\mat}[1]{\mathbf{#1}}
\newcommand{\est}[1]{\widehat{#1}}
\newcommand{\trans}[1]{{#1}^\mathrm{T}}

\newcommand{\sectionfont}{\sffamily\bfseries}
\newcommand{\headerfont}{\sffamily\small}
\newcommand{\footerfont}{\sffamily\small}
\newcommand{\captionfont}{sf,small}


\makeatletter
\expandafter\ifstrequal\expandafter{\myClass}{book}{
    \let\old@frontmatter\frontmatter
    \renewcommand{\frontmatter}{
        \old@frontmatter
            \pagenumbering{roman}
    }
    \let\old@mainmatter\mainmatter
    \renewcommand{\mainmatter}{
            \old@mainmatter
    }
}{}
\makeatother

\usepackage[table]{xcolor}
%main kleuren
\definecolor{purple}{HTML}{8139c0}
%secundaire kleuren
\definecolor{darkblue}{HTML}{003340}
\definecolor{lightpurple}{HTML}{eee4f7}
\definecolor{white}{HTML}{ffffff}
\definecolor{darkpurple}{HTML}{3c1b59}
%accentkleuren
\definecolor{yellow}{HTML}{fcc200}
\definecolor{green}{HTML}{3ab7b0}
\definecolor{blue}{HTML}{00b0eb}

%\usepackage{hyperref}
%\usepackage{cleveref}

\usepackage{tikz}
\usetikzlibrary{shapes, arrows, arrows.meta, positioning, fadings, shadows}
\usepackage{pgfplots}
\pgfplotsset{compat=newest}
\usepackage[american, european resistors, cute inductor, siunitx, nooldvoltagedirection]{circuitikz}
\renewcommand{\familydefault}{\sfdefault}

\usepackage{graphicx} % Required for inserting images
\usepackage{multirow}

\setlength{\parindent}{0 pt}

\pagestyle{empty}

%maken van een titelpagina

%versiebeheer
\newcommand{\versiehistorieheader}{%
	\arrayrulecolor{purple}\hline
	\rowcolor{purple} 
	\multicolumn{1}{c}{\textcolor{white}{\textbf{Datum}}} & \multicolumn{1}{c}{\textcolor{white}{\textbf{Versie}}} & \multicolumn{1}{c}{\textcolor{white}{\textbf{Omschrijving}}} & \multicolumn{1}{c}{\textcolor{white}{\textbf{Auteur}}}
	\\ \hline
}
\newcommand{\versiehistorie}[1]{
	\section*{Versiehistorie}
	{\renewcommand{\arraystretch}{1.24}%
	\vspace*{.1\baselineskip}%
	\begin{xltabular}{\textwidth}{|cc>{\raggedright\arraybackslash}Xc|}
		\versiehistorieheader
		\endfirsthead
		\multicolumn{4}{l}{\hspace*{-\tabcolsep}\sffamily\small\ Vervolg~van~de~vorige~pagina.}\\ \hline
		\versiehistorieheader
		\endhead
		\multicolumn{4}{r}{\sffamily\small Wordt~vervolgd~op~de~volgende~pagina.\hspace*{-\tabcolsep}}\\
		\endfoot
		\hline
		\endlastfoot
		#1%
	\end{xltabular}
	\addtocounter{table}{-1}%
	}\vspace{-.45\baselineskip}%
}

%figuur instellingen
% iets meer ruimte boven en iets minder onder figuur:
\makeatletter
\newenvironment{myFigure}[1][!htbp]{%
	\expandafter\figure\expandafter[#1]%
    \expandafter\ifstrequal\expandafter{#1}{H}{%
        \vspace{.2\baselineskip}%
    }{%
		\vspace{.3\baselineskip}%
    }%
	\def\end@fig@space{%
        \expandafter\ifstrequal\expandafter{#1}{H}{%
            \vspace{-.85\baselineskip}%
        }{%
            \vspace{-.2\baselineskip}%
        }%
    }%
}{%
    \end@fig@space{}%
	\endfigure%
}
\makeatother

%Circuit instellingen
% iets meer ruimte boven en iets minder onder figuur:
\makeatletter
\newenvironment{myCircuit}[1][!htbp]{%
	\expandafter\circuit\expandafter[#1]%
    \expandafter\ifstrequal\expandafter{#1}{H}{%
        \vspace{.2\baselineskip}%
    }{%
		\vspace{.3\baselineskip}%
    }%
	\def\end@fig@space{%
        \expandafter\ifstrequal\expandafter{#1}{H}{%
            \vspace{-.85\baselineskip}%
        }{%
            \vspace{-.2\baselineskip}%
        }%
    }%
}{%
    \end@fig@space{}%
	\endcircuit%
}
\makeatother

%%>tabel instellingen
% iets meer ruimte boven en iets minder onder tabel:
\makeatletter
\newenvironment{myTable}[1][!htbp]{%
    \begin{table}[#1]%
    \expandafter\ifstrequal\expandafter{#1}{H}{%
	    \vspace{.1\baselineskip}%
	}{%
    	\vspace{.2\baselineskip}%
	}%
    \def\end@table@space{%
        \expandafter\ifstrequal\expandafter{#1}{H}{%
            \vspace{-.6\baselineskip}%
        }{%
            \vspace{.1\baselineskip}%
        }%
    }%
}{%
    \end@table@space{}%
    \end{table}%
}
\makeatother

%>figuur
% commando voor figuur
% figuren staan in directory ./figs
% #1 OPTIONEEL = placement argumenten voor figure b.v. tb of H 
% #2 = width of scale optie voor \includegraphics
% #3 = filenaam (zonder extensie); labelnaam=fig:filenaam
% #4 = titel
\newcommand{\figuur}[4][!htbp]{
    \begin{myFigure}[#1]
        \centering%
        \includegraphics[#2]{figs/#3}
        \caption{#4}
        \label{fig:#3}
    \end{myFigure}
}
%<figuur

%>tikzfiguur
% commando voor figuur getekend met tikz
% figuren staan in directory ./tikz
% #1 OPTIONEEL = placement argumenten voor figure b.v. tb of H 
% #2 = filenaam (zonder extensie); labelnaam=fig:filenaam
% #3 = titel
\newcommand{\tikzfiguur}[3][!htbp]{%
	\begin{myFigure}[#1]%
		\centering%
		\input{tikz/#2}%
		\caption{#3}%
		\label{fig:#2}%
	\end{myFigure}%
}
%<tikzfiguur

%>tikzinlinefiguur
% commando voor figuur getekend met tikz
% figuren staan inline
% #1 OPTIONEEL = placement argumenten voor figure b.v. tb of H 
% #2 = label
% #3 = TikZ-code
% #4 = titel
\newcommand{\tikzinlinefiguur}[4][H]{
	\begin{myFigure}[#1]%
		\centering%
		\begin{tikzpicture}%
		#3%
		\end{tikzpicture}%
		\caption{#4}%
		\label{fig:#2}%
	\end{myFigure}
}
%<tikzinlinefiguur

%>schakeling
% commando voor schakeling
% #1 OPTIONEEL = placement argumenten voor figure b.v. tb of H 
% #2 = circuittikz code
% #3 = titel
% #4 = labelnaam=sch:#4
% spacing voor en na schakeling consistent gemaakt met figuur
\newcommand{\schakeling}[4][!htbp]{
	\begin{myCircuit}[#1]%
        \centering%
		\begin{circuitikz}%
			#2%
		\end{circuitikz}%	
		\caption{#3}%
		\label{sch:#4}%
	\end{myCircuit}%
}
\makeatother
%<schakeling

\usepackage{tabularx}
\usepackage{xltabular}

\newcommand{\USE}[1]{\foreignlanguage{USenglish}{#1}}
\newcommand{\hyp}{-\penalty0\hskip0pt\relax}

\newcommand{\csecref}[1]{%
    \cref{#1}%
    \iftoggle{ebook}{}{
        op \cpageref{#1}%
    }%
}

\newcommand{\Csecref}[1]{%
    \Cref{#1}%
    \iftoggle{ebook}{}{
        op \cpageref{#1}%
    }%
}

%>pdfpagref
% verwijzingen naar specifieke pagina van een ander pdf document op internet
% #1 = OPTIONEEL Verschil tussen paginanummer (weergegeven in pdf viewer) en nummer weergegeven op de pagina zelf.
% #2 = url van het betreffende pdf document
% #3 = nummer weergegeven op de pagina zelf
\makeatletter
\newcounter{pdfpagref@counter}
\newcommand{\pdfpagref}[3][0]{%
    \setcounter{pdfpagref@counter}{#3}%
    \addtocounter{pdfpagref@counter}{#1}%
    \href{#2\#page=\thepdfpagref@counter}{pagina~#3}%
}
\makeatother
%<pdfpagref

%>pdfdeelref
% verwijzingen naar specifiek deel van een ander pdf document op internet
% #1 = naam van deel, b.v. chapter, section, subsection, enz
% #2 = naam voor deel,  b.v. hoofdstuk~, paragraaf~, enz
% #3 = url van het betreffende pdf document
% #4 = nummer van het betreffende deel, bijvoorbeeld 6.2
\newcommand{\pdfdeelref}[4]{%
    \href{#3\#nameddest=#1.#4}{#2#4}%
}
%<pdfdeelref

\usepackage[USenglish, dutch]{babel}
\usepackage[
    backend=biber,
    backref=true,
    backrefstyle=none, % neem paginanummer niet samen in de backref lijst dus 1,2,3 i.p.v. 1-3 (anders kun je niet op 2 klikken)
    sortcites=true,
    doi = false, % doi informatie wordt niet weergegeven
    useprefix=false
]{biblatex}
\addbibresource{bibliography.bib}
\makeatletter
\AtBeginDocument{\toggletrue{blx@useprefix}}
\makeatother
\DefineBibliographyStrings{dutch}{
    backrefpage = {geciteerd op p.},
    backrefpages = {geciteerd op pp.},
}

%>citepdfpag
% verwijzingen naar specifieke pagina van
% een ander pdf document op internet uit de referentielijst
% #1 = OPTIONEEL Verschil tussen paginanummer (weergegeven in pdf viewer) en nummer weergegeven op de pagina zelf.
% #2 = Bibtexkey
% #3 = nummer weergegeven op de pagina zelf
% geinspireerd door: http://tex.stackexchange.com/questions/191174/biblatex-citefield-not-expanded-as-expected
\makeatletter
\DeclareCiteCommand{\@saveurlcite}{}{\savefield{url}{\@savedurl}}{}{}
\newcounter{cite@pag}
\newcommand{\citepdfpag}[3][0]{%
    \setcounter{cite@pag}{#3}%
    \addtocounter{cite@pag}{#1}%
    \@saveurlcite{#2}%
    \edef\@savedurlpag{\@savedurl\#page=\thecite@pag}%
    \cite[\href{\@savedurlpag}{pagina~#3}]{#2}%
}
\makeatother
%<citepdfpag

%>citepdfdeel
% verwijzingen naar specifiek deel van
% een ander pdf document op internet uit de referentielijst
% #1 = naam van deel, b.v. chapter, section, subsection, enz
% #2 = naam voor deel,  b.v. hoofdstuk~, paragraaf~, enz
% #3 = Bibtexkey van het betreffende pdf document
% #4 = nummer van het betreffende deel, bijvoorbeeld 6.2
% geinspireerd door: http://www.tex.ac.uk/FAQ-twooptarg.html
\makeatletter
\newcommand\citepdfdeel[4]{%
    \@saveurlcite{#3}%
    \edef\@savedurldeel{\@savedurl\#nameddest=#1.#4}%
    \cite[\href{\@savedurldeel}{#2#4}]{#3}%
}
\makeatother
%<citepdfdeel

\usepackage{titlesec}
\titleformat{\chapter}[display]{\sectionfont\Huge}{%
\sectionfont\fontsize{45.64172pt}{0pt}\selectfont\filleft\thechapter}{.5\baselineskip}{}
{}
\titleformat{\section}{\sectionfont\Large}{\thesection}{1em}{}
\titleformat{\subsection}{\sectionfont\large}{\thesubsection}{1em}{}
\titleformat{\subsubsection}{\sectionfont\large}{\thesubsubsection}{1em}{}
%mogelijk anpassing opmaak titels

\makeatletter
\begingroup
\newif\ifmiktex
\def\MiKTeX{MiKTeX}
\@onelevel@sanitize\MiKTeX
% Scans for 'MiKTeX' string with catcode 12 (other)
\expandafter\def\expandafter\testmiktex\expandafter#\expandafter1\MiKTeX#2\relax{%
    \ifx\empty#2\empty
        \miktexfalse
    \else
        \miktextrue
    \fi
}
\expandafter\expandafter
\expandafter\testmiktex\expandafter\pdftexbanner\MiKTeX\relax\relax
\ifmiktex
    \gdef\TeXLiveVersion{0}
\else
    \edef\get@@version{\def\noexpand\get@version##1\detokenize{TeX Live}##220##3##4##5)##6\noexpand\@nil}
    \get@@version{\gdef\TeXLiveVersion{20#3#4}}
    \expandafter\get@version\pdftexbanner\@nil
\fi
\endgroup
\makeatother

% ruimte onder aan de pagina tussen tekst en voetnoot niet na voetnoot
% package moet VOOR fancyhdr om warning te voorkomen
\usepackage[
    bottom,
    hang,
    multiple
]{footmisc}
% inspringen
\setlength{\footnotemargin}{1em}
% ruimte tussen footnotes:
\setlength{\footnotesep}{.8\baselineskip}

\expandafter\ifstrequal\expandafter{\myClass}{book}{
% Voetnoot teller NIET resetten bij elk hoofdstuk
\let\counterwithout\relax
\let\counterwithin\relax
\usepackage{chngcntr}
\counterwithout{footnote}{chapter}
}


\makeatletter

\newlength{\page@number@width}
\settowidth{\page@number@width}{\thepage}
\addtolength{\page@number@width}{6mm}

\newcommand{\rease@header}{\vspace{4pt}}

\newcommand{\page@number@box}{\colorbox{purple}{\makebox[\page@number@width][c]{\color{white}\textbf{\vphantom{2}\thepage}}}}

\usepackage{fancyhdr}
\newcommand{\myfancy}{
    \fancyhead{}
    \fancyfoot{}
		\expandafter\ifstrequal\expandafter{\myClass}{article}{
        	\fancyhead[R]{\headerfont\nouppercase{\rightmark}\rease@header{}}
		}{
        	\fancyhead[R]{\headerfont\nouppercase{\leftmark}\rease@header{}}
		}
        \fancyfoot[L]{\headerfont{}RMI latex} %bedrijf/ vak ipv school
        \fancyfoot[R]{\headerfont\titel{}\hspace{1em}\page@number@box{}}
    \renewcommand{\headrulewidth}{0.4pt}
    \renewcommand{\footrulewidth}{0.4pt}
}
\makeatother
\markboth{\text{}}{\text{}}
\pagestyle{fancy}
\myfancy
\fancypagestyle{plain}{\myfancy}

\usepackage{parskip}
\usepackage{graphicx}
\usepackage{array}
\usepackage{longtable}
\usepackage{multirow}
\usepackage{adjustbox}
\usepackage{float}

\renewcommand{\topfraction}{.95} %in plaats van .70
\renewcommand{\bottomfraction}{.95} %in plaats van .30
\renewcommand{\textfraction}{.05} %in plaats van .20
\renewcommand{\floatpagefraction}{.90} %in plaats van .5
\usepackage{siunitx}

\sisetup{
    per-mode = symbol,
% binary-units = true,
    range-units = single,
% dutch localisation
    range-phrase = { tot en met },
    output-decimal-marker = {,},
    list-final-separator = { en },
    list-pair-separator  = { en },
}

\usepackage[
    font={\captionfont},
    format=hang,
    justification=RaggedRight,
    labelfont={color=purple, bf},
    hypcap=true,
    skip=.6\baselineskip
]{caption}

% vermijdt weduwen en wezen
\usepackage[defaultlines=2,all]{nowidow}

% standaard Nederlandse quotes zijn old-fashioned (openen laag)
%>csquote
\usepackage[style=english]{csquotes}
\newcommand{\dquote}[1]{\enquote{#1}}
\newcommand{\squote}[1]{\enquote*{#1}}
%<csquote

% enumitem zorgt ervoor dat je een enumeratie eenvoudiger kunt aanpassen. B.v. het inspringen.
\usepackage{enumitem}
\setlistdepth{5}
\setlist[enumerate,1]{label=\arabic*.}
\setlist[enumerate,2]{label=\arabic*.}
\setlist[enumerate,3]{label=\arabic*.}
\setlist[enumerate,4]{label=\arabic*.}
\setlist[enumerate,5]{label=\arabic*.}
\renewlist{enumerate}{enumerate}{5}
\setlist{topsep=-.2\baselineskip, itemsep=-.3\baselineskip}
\setlist[itemize,1]{label=\textbullet}
\setlist[itemize,2]{label=$\circ$}
\setlist[itemize,3]{label=$\triangleright$}
\setlist[itemize,4]{label=$\diamond$}

\usepackage{footnotebackref}
% patch see:
\makeatletter
\LetLtxMacro{\BHFN@Old@footnotemark}{\@footnotemark}
\renewcommand*{\@footnotemark}{%
    \refstepcounter{BackrefHyperFootnoteCounter}%
    \xdef\BackrefFootnoteTag{bhfn:\theBackrefHyperFootnoteCounter}%
    \label{\BackrefFootnoteTag}%
    \BHFN@Old@footnotemark
}
\makeatother

\usepackage{microtype}
\usepackage{hyperref}

\hypersetup{
    colorlinks=true,
    breaklinks=true,
    anchorcolor=blue, 
    citecolor=blue,
    filecolor=blue,
    linkcolor=blue,
    urlcolor=blue,
    pdftitle={\titel{}},
    pdfauthor={Hogeschool Rotterdam},
    pdfdisplaydoctitle=true,
    pdfpagelayout=\iftoggle{onesided}{SinglePage}{TwoPageRight},
    %pdfstartview={XYZ null null 1.50}, %voor debuggen
    %pdfstartview=FitH, % voor debuggen
    }

    \pdfsuppressptexinfo=1

    % url met mogelijkheid tot afbreken:
\newcommand{\urlhyp}[1]{%
	\StrSubstitute{#1}{\\-}{}[\nakedurl]%
	\href{\nakedurl}{\hcode{#1}}%
}

% href met voetnoot met daarin de URL:
\newcommand{\hreffn}[3][]{%
	\StrSubstitute{#2}{\\-}{}[\nakedurl]%
	\href[#1]{#2}{#3}\footnote{\href{\nakedurl}{\hcode{#2}}}%
}

\usepackage[dutch, noabbrev, nameinlink]{cleveref}
\crefname{lstlisting}{listing}{listings}
\Crefname{lstlisting}{Listing}{Listings}
\crefname{opdracht}{opdracht}{opdrachten}
\Crefname{opdracht}{Opdracht}{Opdrachten}
\expandafter\ifstrequal\expandafter{\myClass}{article}{
    \crefname{section}{hoofdstuk}{hoofdstukken}
    \Crefname{section}{Hoofdstuk}{Hoofdstukken}
    \crefname{subsection}{paragraaf}{paragrafen}
    \Crefname{subsection}{Paragraaf}{Paragrafen}
    \crefname{subsubsection}{paragraaf}{paragrafen}
    \Crefname{subsubsection}{Paragraaf}{Paragrafen}
}{}
\crefname{appendix}{bijlage}{bijlagen}
\Crefname{appendix}{Bijlage}{Bijlagen}
\crefname{circuit}{schakeling}{schakelingen}
\Crefname{circuit}{Schakeling}{Schakelingen}
\newcommand{\crefrangeconjunction}{ tot en met }